\chapter{Способы и методы теоретических и аналитических исследований} \label{ch:methods}
% (для теоретических и аналитических работ)
Раздел должен содержать обоснование и подробную пошаговую реализацию
примененного теоретического и/или аналитического метода. Описание метода должно
обеспечивать возможность независимого воспроизведения результатов, полученных в
работе.

% \chapter{Методика эксперимента} \label{ch:methods}
% % (для экспериментальных, аналитических и теоретических работ, в том случае,
% % если в них проводится натурный и/или численный эксперимент)
% Методика должна описывать постановку натурного и/или численного эксперимента,
% описание и схемы оригинальных приборов и установок, тип и модели
% использованного стандартного оборудования, порядок проведения измерений и
% обработки результатов, анализ погрешностей натурного и/или численного
% эксперимента. Описание методики эксперимента должно обеспечивать возможность
% независимого воспроизведения экспериментальных данных, полученных в работе.

% \chapter{Методика проведения наблюдений и измерений} \label{ch:methods}
% % (для экспериментальных, аналитических и теоретических работ, в том случае,
% % если в них проводятся наблюдения естественных и искусственных объектов и
% % измерение их параметров)
% Методика должна включать в себя описание и схемы оригинальных приборов и
% установок, тип и модели использованного стандартного оборудования, порядок
% проведения наблюдений, измерений и обработки результатов, анализ погрешностей
% наблюдений и измерений. Описание методики наблюдений и измерений должно
% обеспечивать возможность независимого воспроизведения экспериментальных данных,
% полученных в работе.

% \chapter{Способы и методы решения задачи} \label{ch:methods}
% % (для экспериментальных и инженерно-практических задач)
% Раздел должен содержать обоснование и подробную пошаговую реализацию
% примененного метода решения задачи квалификационной работы. Описание метода
% должно обеспечивать возможность независимого воспроизведения результатов,
% полученных в работе.

